% Options for packages loaded elsewhere
\PassOptionsToPackage{unicode}{hyperref}
\PassOptionsToPackage{hyphens}{url}
\PassOptionsToPackage{dvipsnames,svgnames,x11names}{xcolor}
%
\documentclass[
  10pt,
]{article}
\usepackage{amsmath,amssymb}
\usepackage{iftex}
\ifPDFTeX
  \usepackage[T1]{fontenc}
  \usepackage[utf8]{inputenc}
  \usepackage{textcomp} % provide euro and other symbols
\else % if luatex or xetex
  \usepackage{unicode-math} % this also loads fontspec
  \defaultfontfeatures{Scale=MatchLowercase}
  \defaultfontfeatures[\rmfamily]{Ligatures=TeX,Scale=1}
\fi
\usepackage{lmodern}
\ifPDFTeX\else
  % xetex/luatex font selection
\fi
% Use upquote if available, for straight quotes in verbatim environments
\IfFileExists{upquote.sty}{\usepackage{upquote}}{}
\IfFileExists{microtype.sty}{% use microtype if available
  \usepackage[]{microtype}
  \UseMicrotypeSet[protrusion]{basicmath} % disable protrusion for tt fonts
}{}
\makeatletter
\@ifundefined{KOMAClassName}{% if non-KOMA class
  \IfFileExists{parskip.sty}{%
    \usepackage{parskip}
  }{% else
    \setlength{\parindent}{0pt}
    \setlength{\parskip}{6pt plus 2pt minus 1pt}}
}{% if KOMA class
  \KOMAoptions{parskip=half}}
\makeatother
\usepackage{xcolor}
\usepackage[margin=0.78in]{geometry}
\usepackage{color}
\usepackage{fancyvrb}
\newcommand{\VerbBar}{|}
\newcommand{\VERB}{\Verb[commandchars=\\\{\}]}
\DefineVerbatimEnvironment{Highlighting}{Verbatim}{commandchars=\\\{\}}
% Add ',fontsize=\small' for more characters per line
\newenvironment{Shaded}{}{}
\newcommand{\AlertTok}[1]{\textcolor[rgb]{1.00,0.00,0.00}{\textbf{#1}}}
\newcommand{\AnnotationTok}[1]{\textcolor[rgb]{0.38,0.63,0.69}{\textbf{\textit{#1}}}}
\newcommand{\AttributeTok}[1]{\textcolor[rgb]{0.49,0.56,0.16}{#1}}
\newcommand{\BaseNTok}[1]{\textcolor[rgb]{0.25,0.63,0.44}{#1}}
\newcommand{\BuiltInTok}[1]{\textcolor[rgb]{0.00,0.50,0.00}{#1}}
\newcommand{\CharTok}[1]{\textcolor[rgb]{0.25,0.44,0.63}{#1}}
\newcommand{\CommentTok}[1]{\textcolor[rgb]{0.38,0.63,0.69}{\textit{#1}}}
\newcommand{\CommentVarTok}[1]{\textcolor[rgb]{0.38,0.63,0.69}{\textbf{\textit{#1}}}}
\newcommand{\ConstantTok}[1]{\textcolor[rgb]{0.53,0.00,0.00}{#1}}
\newcommand{\ControlFlowTok}[1]{\textcolor[rgb]{0.00,0.44,0.13}{\textbf{#1}}}
\newcommand{\DataTypeTok}[1]{\textcolor[rgb]{0.56,0.13,0.00}{#1}}
\newcommand{\DecValTok}[1]{\textcolor[rgb]{0.25,0.63,0.44}{#1}}
\newcommand{\DocumentationTok}[1]{\textcolor[rgb]{0.73,0.13,0.13}{\textit{#1}}}
\newcommand{\ErrorTok}[1]{\textcolor[rgb]{1.00,0.00,0.00}{\textbf{#1}}}
\newcommand{\ExtensionTok}[1]{#1}
\newcommand{\FloatTok}[1]{\textcolor[rgb]{0.25,0.63,0.44}{#1}}
\newcommand{\FunctionTok}[1]{\textcolor[rgb]{0.02,0.16,0.49}{#1}}
\newcommand{\ImportTok}[1]{\textcolor[rgb]{0.00,0.50,0.00}{\textbf{#1}}}
\newcommand{\InformationTok}[1]{\textcolor[rgb]{0.38,0.63,0.69}{\textbf{\textit{#1}}}}
\newcommand{\KeywordTok}[1]{\textcolor[rgb]{0.00,0.44,0.13}{\textbf{#1}}}
\newcommand{\NormalTok}[1]{#1}
\newcommand{\OperatorTok}[1]{\textcolor[rgb]{0.40,0.40,0.40}{#1}}
\newcommand{\OtherTok}[1]{\textcolor[rgb]{0.00,0.44,0.13}{#1}}
\newcommand{\PreprocessorTok}[1]{\textcolor[rgb]{0.74,0.48,0.00}{#1}}
\newcommand{\RegionMarkerTok}[1]{#1}
\newcommand{\SpecialCharTok}[1]{\textcolor[rgb]{0.25,0.44,0.63}{#1}}
\newcommand{\SpecialStringTok}[1]{\textcolor[rgb]{0.73,0.40,0.53}{#1}}
\newcommand{\StringTok}[1]{\textcolor[rgb]{0.25,0.44,0.63}{#1}}
\newcommand{\VariableTok}[1]{\textcolor[rgb]{0.10,0.09,0.49}{#1}}
\newcommand{\VerbatimStringTok}[1]{\textcolor[rgb]{0.25,0.44,0.63}{#1}}
\newcommand{\WarningTok}[1]{\textcolor[rgb]{0.38,0.63,0.69}{\textbf{\textit{#1}}}}
\usepackage{longtable,booktabs,array}
\usepackage{calc} % for calculating minipage widths
% Correct order of tables after \paragraph or \subparagraph
\usepackage{etoolbox}
\makeatletter
\patchcmd\longtable{\par}{\if@noskipsec\mbox{}\fi\par}{}{}
\makeatother
% Allow footnotes in longtable head/foot
\IfFileExists{footnotehyper.sty}{\usepackage{footnotehyper}}{\usepackage{footnote}}
\makesavenoteenv{longtable}
\setlength{\emergencystretch}{3em} % prevent overfull lines
\providecommand{\tightlist}{%
  \setlength{\itemsep}{0pt}\setlength{\parskip}{0pt}}
\setcounter{secnumdepth}{5}
\usepackage{booktabs}
\usepackage{longtable}
\usepackage{array}
\usepackage{microtype}
\usepackage{enumitem}
\setlist{nosep,leftmargin=*}
\setlength{\parindent}{0pt}
\setlength{\parskip}{0.45em}
\ifLuaTeX
  \usepackage{selnolig}  % disable illegal ligatures
\fi
\usepackage{bookmark}
\IfFileExists{xurl.sty}{\usepackage{xurl}}{} % add URL line breaks if available
\urlstyle{same}
\hypersetup{
  pdftitle={SRIC Block 24},
  pdfauthor={Komorebi P.},
  colorlinks=true,
  linkcolor={blue},
  filecolor={Maroon},
  citecolor={Blue},
  urlcolor={blue},
  pdfcreator={LaTeX via pandoc}}

\title{SRIC Block 24}
\usepackage{etoolbox}
\makeatletter
\providecommand{\subtitle}[1]{% add subtitle to \maketitle
  \apptocmd{\@title}{\par {\large #1 \par}}{}{}
}
\makeatother
\subtitle{External Latent-Reasoner Causal-Semantic Admission - Protocol
v0.3.0}
\author{Komorebi P.}
\date{September 3, 2026}

\begin{document}
\maketitle

{
\hypersetup{linkcolor=}
\setcounter{tocdepth}{3}
\tableofcontents
}
\section{Executive status}\label{executive-status}

\textbf{Canonical status:}
\texttt{INFRASTRUCTURE\_VALIDATED\_EXTERNAL\_EXECUTION\_PENDING}

Block 24 is the first SRIC package whose intended evidence source is an
independently authored and independently trained latent-reasoning
checkpoint. It translates the synthetic Neuralese program of Blocks
16-23 into a source-locked external-model admission protocol.

The public checkpoint was \textbf{not executed in the build
environment}. The release therefore reports no empirical result about
Coconut, PrOntoQA, or any other third-party latent reasoner. It
establishes the external adapter, source locks, frozen chronology,
causal intervention battery, semantic escalation, post-seal witnesses,
engineered positive control, and reproducibility machinery required for
that run.

A negative real result is fully admissible. The pipeline is designed to
classify failure rather than force a positive interpretation.

\section{Research strategy}\label{research-strategy}

The next major gap is external validity. Another package-local positive
control would contribute less evidence than exposing the architecture to
an independently authored recurrence and checkpoint. The design
therefore prioritizes source fidelity and causal clarity over model
size.

The conservative defaults are:

\begin{enumerate}
\def\labelenumi{\arabic{enumi}.}
\tightlist
\item
  GPT-2 multimode Coconut on PrOntoQA is the primary target.
\item
  PrOntoQA supports a \texttt{CONTENT\_SPECIFIC} claim, not a
  different-algorithm claim.
\item
  Explicit-CoT is revealed only after the semantic quotient is sealed.
\item
  ProsQA is the first structure-sensitive \texttt{PROCEDURE\_SPECIFIC}
  follow-up.
\item
  Smoke, primary, and full profiles are run in that order without moving
  scientific thresholds.
\end{enumerate}

\section{Part I - Brainstorm and design
decisions}\label{part-i---brainstorm-and-design-decisions}

\subsection{Why the next gap is
empirical}\label{why-the-next-gap-is-empirical}

Blocks 16-23 progressively established a synthetic theory of latent
causal auditing, semantic charts, cross-architecture transport,
obstruction-guided repair, unpaired alignment, and endogenous semantic
quotient construction. Another synthetic positive control would add less
evidence than testing whether the same gates survive contact with a
checkpoint designed by an independent group.

The central next question is therefore not whether another controlled
fixture can be made to pass. It is whether a source-locked external
latent reasoner lands in one of SRIC's predeclared epistemic classes
without altering the gates after seeing the result.

\subsection{Candidate external models
considered}\label{candidate-external-models-considered}

\subsubsection{Multimode GPT-2 Coconut on
PrOntoQA}\label{multimode-gpt-2-coconut-on-prontoqa}

Advantages:

\begin{itemize}
\tightlist
\item
  small enough for an 8 GB laptop GPU at batch size one;
\item
  public repository and raw checkpoint;
\item
  latent, direct-answer, and explicit-CoT modes in one trained wrapper;
\item
  recurrent latent vectors are reinjected at identifiable latent-token
  positions;
\item
  logic answers and proof chains are externally checkable;
\item
  semantic twins and counterfactual query interventions are easy to
  construct.
\end{itemize}

Limitation:

\begin{itemize}
\tightlist
\item
  the task family is dominated by fixed-form unary chains. It supports a
  clean first causal audit but not a strong claim that different latent
  traces encode different reasoning algorithms.
\end{itemize}

Decision: use this checkpoint as the primary source-faithful causal
admission target, with maximum causal status \texttt{CONTENT\_SPECIFIC}
before semantic escalation.

\subsubsection{ProsQA multimode Coconut}\label{prosqa-multimode-coconut}

Advantages:

\begin{itemize}
\tightlist
\item
  more varied proof structures;
\item
  better setting for structure-near versus structure-far donor
  contrasts;
\item
  same external research codebase and model family.
\end{itemize}

Limitations:

\begin{itemize}
\tightlist
\item
  answer parsing and proof-structure extraction require a separate
  adapter path;
\item
  harder to debug simultaneously with the first source-faithful
  recurrence adapter.
\end{itemize}

Decision: reserve as Tier 2. It is the first dataset on which
\texttt{PROCEDURE\_SPECIFIC} should be enabled.

\subsubsection{SLPO Coconut}\label{slpo-coconut}

Advantages:

\begin{itemize}
\tightlist
\item
  exposes adaptive stopping information;
\item
  provides latent lengths and gate outputs.
\end{itemize}

Limitations:

\begin{itemize}
\tightlist
\item
  additional runtime and checkpoint family;
\item
  adaptive stopping introduces another mechanism before basic replay is
  proven;
\item
  full training experiments exceed the current 8 GB hardware target.
\end{itemize}

Decision: use after the fixed-length causal audit is stable.

\subsubsection{Hyper 3B}\label{hyper-3b}

Advantages:

\begin{itemize}
\tightlist
\item
  directly recurrent latent reasoning;
\item
  larger and more ambitious external stress.
\end{itemize}

Limitations:

\begin{itemize}
\tightlist
\item
  at the hardware boundary;
\item
  adapter and memory failures could be confused with theoretical
  failure;
\item
  current halt behavior is not a strong adaptive-depth positive control.
\end{itemize}

Decision: later stress test, not the first empirical target.

\subsection{Major methodological
choices}\label{major-methodological-choices}

\subsubsection{1. Adapter fidelity before model
ambition}\label{adapter-fidelity-before-model-ambition}

A large checkpoint is scientifically useless if the intervention adapter
does not reproduce the model's native recurrence. Block 24 therefore
gives exact replay and source-code parity priority over scale.

\subsubsection{2. Teacher-forced answer loss as the primary effect
measure}\label{teacher-forced-answer-loss-as-the-primary-effect-measure}

Free generation is useful descriptively but introduces parsing,
decoding, and stochasticity. The main causal effects use teacher-forced
protected answer loss while holding the prompt and suffix constant.
Free-generation accuracy is a secondary behavioral witness.

\subsubsection{3. Content specificity is not procedure
specificity}\label{content-specificity-is-not-procedure-specificity}

A same-answer donor rules out pure final-answer encoding. It does not
automatically prove a different algorithm is represented. On PrOntoQA
the result is called content/proof-instance specificity. Procedure
specificity is reserved for a task family with varied proof topology and
an explicit far-versus-near structural contrast.

\subsubsection{4. Explicit CoT is post-seal
only}\label{explicit-cot-is-post-seal-only}

The multimode checkpoint can be given gold proof steps in verbalized
mode. Allowing those traces into semantic construction would leak the
target ontology. Therefore explicit CoT is revealed only after the
quotient seal as a comparison witness.

\subsubsection{5. Covariance calibration is
disjoint}\label{covariance-calibration-is-disjoint}

Covariance-matched noise must not estimate its nuisance distribution
from the same rows used to measure the final effect. Twenty percent of
proof groups are deterministically reserved for covariance calibration.

\subsubsection{6. Semantic final traces are not computed before
sealing}\label{semantic-final-traces-are-not-computed-before-sealing}

The previous prototype split the records correctly but captured every
semantic orbit before choosing the quotient. The v0.3 implementation now
captures discovery, then calibration, seals the selected quotient, and
only then captures final semantic twins and counterfactuals.

\subsubsection{7. Fail closed}\label{fail-closed}

Replay failure terminates the causal claim. Missing source objects
terminate setup. Dirty or wrong upstream commits terminate source
verification. No semantic candidate produces
\texttt{ONTOLOGY\_UNDERIDENTIFIED}, not a threshold adjustment.

\subsubsection{8. Smallest valid
quotient}\label{smallest-valid-quotient}

Semantic construction is validity first, then minimum rank. Full hidden
state is a descriptive upper control and is ineligible for selection.

\subsection{Alternatives rejected for the first
run}\label{alternatives-rejected-for-the-first-run}

\begin{itemize}
\tightlist
\item
  Training a new latent model: confounds model construction with
  external validation.
\item
  Beginning with a 3B checkpoint: raises adapter and memory risk without
  improving causal clarity.
\item
  Using explicit CoT during quotient discovery: leaks proof semantics.
\item
  Calling any same-answer transplant procedure-specific: overstates what
  fixed-chain tasks identify.
\item
  Estimating covariance from the final audit: weakens the intervention's
  independence.
\item
  Selecting only tasks where latent recurrence helps: creates
  post-selection bias. The full causal population is reported first.
\item
  Lowering gates after failure: invalidates the sealed experiment.
\end{itemize}

\subsection{Remaining judgment calls}\label{remaining-judgment-calls}

The user was asked whether to prioritize a small clean causal proof or
the strongest latent model. No answer was available before
implementation, so the conservative default was chosen: small and
source-faithful first.

The user was also asked whether explicit CoT should help construct
semantics or remain post-seal. The conservative default is post-seal
only.

The next strategic choice after the PrOntoQA primary result is whether
to proceed immediately to ProsQA or first expand the PrOntoQA final
population. The recommended default is to inspect the primary receipt,
then use ProsQA only if replay and at least useful ordered recurrence
are supported.

\newpage

\section{Part II - Canonical
protocol}\label{part-ii---canonical-protocol}

\subsection{External Latent-Reasoner Causal-Semantic
Admission}\label{external-latent-reasoner-causal-semantic-admission}

\subsubsection{1. Purpose}\label{purpose}

Block 24 is the first SRIC protocol designed to admit or reject a
third-party latent-reasoning checkpoint under a frozen causal and
semantic standard. The primary target is the multimode GPT-2 Coconut
model trained on PrOntoQA. The protocol treats a negative result as
complete evidence when the failure is reproducible and correctly
localized.

\subsubsection{2. Primary scientific
question}\label{primary-scientific-question}

Does the recurrent latent trace of the external checkpoint carry
content-specific, ordered, structured information that causally
contributes to its answer and supports a post-seal low-rank semantic
quotient?

The protocol does not initially ask whether every latent step has a
unique English translation. It asks whether the latent computation has
operationally identifiable causal content.

\subsubsection{3. Source lock}\label{source-lock}

The external model is admitted only when all of the following match the
bundled lock:

\begin{itemize}
\tightlist
\item
  repository: \texttt{connordilgren/are-lrms-easily-interpretable};
\item
  commit: \texttt{32f413d8d55239d9bc54bb6b6ec37b0630891ed4};
\item
  multimode wrapper Git blob:
  \texttt{d0b353cd78f98ab7aae0d20ce4e8859d9fd9958e};
\item
  Coconut recurrence Git blob:
  \texttt{646acb08f149afec075eefb212c8aa1f3ac8ceda};
\item
  requirements Git blob:
  \texttt{4a63ea4b92f0dd565f5d632241e36299695237e4};
\item
  PrOntoQA test-data Git blob:
  \texttt{7222832226883a545328f2398dc4976b1bb85c7d};
\item
  checkpoint repository revision:
  \texttt{4ec759a9f9b8a8e7e1a8d1e3e2d43eef969ac83f};
\item
  checkpoint filename: \texttt{checkpoint\_37};
\item
  GPT-2 base revision:
  \texttt{607a30d783dfa663caf39e06633721c8d4cfcd7e}.
\end{itemize}

The downloaded checkpoint is hashed in full. The complete local
base-model payload is hashed as a path-size-file-digest tree with
download metadata excluded.

\subsubsection{4. Adapter contract}\label{adapter-contract}

The adapter must expose the exact vectors reinjected at latent-token
positions. It must support:

\begin{itemize}
\tightlist
\item
  ordinary latent capture;
\item
  exact supplied-trace replay;
\item
  direct-answer mode;
\item
  post-seal verbalized mode;
\item
  teacher-forced protected answer scoring.
\end{itemize}

The trace engine is compared against a separately written reference
implementation of the reviewed source recurrence. Exact replay is then
tested on every audit row. A replay error above the frozen tolerance
returns \texttt{ADAPTER\_INVALID} and terminates all later
interpretation.

\subsubsection{5. Population and
grouping}\label{population-and-grouping}

Tasks are selected deterministically from the pinned public data. Splits
are made by \texttt{group\_id}, not by transformed text row, so semantic
twins and related records cannot cross evidence partitions.

For the primary 64-task profile:

\begin{itemize}
\tightlist
\item
  20\% of groups are reserved for covariance calibration;
\item
  the remainder forms the causal audit population;
\item
  if causal admission succeeds, that population is split into 50\%
  semantic discovery, 25\% semantic calibration, and 25\% semantic final
  support.
\end{itemize}

The exact row counts depend on group sizes and are frozen in the
pre-audit seal.

\subsubsection{6. Causal interventions}\label{causal-interventions}

For task \texttt{i} with clean trace \texttt{H\_i} and clean answer loss
\texttt{L\_i}, the protocol computes:

\paragraph{Exact replay}\label{exact-replay}

\texttt{r\_i\ =\ max\_abs(logits(clean)\ -\ logits(replay(H\_i)))}.

\paragraph{Latent necessity}\label{latent-necessity}

\texttt{n\_i\ =\ L\_i(zero\ trace)\ -\ L\_i}.

\paragraph{Step order}\label{step-order}

\texttt{o\_i\ =\ L\_i(permuted\ trace)\ -\ L\_i}.

\paragraph{Norm-matched structure}\label{norm-matched-structure}

\texttt{c\_i\ =\ L\_i(norm-matched\ noise)\ -\ L\_i}.

\paragraph{Covariance-matched
structure}\label{covariance-matched-structure}

A nuisance covariance is estimated only on the reserved calibration
support. The audit effect is:

\texttt{c\_cov\_i\ =\ L\_i(covariance-matched\ noise)\ -\ L\_i}.

\paragraph{Same-answer content
transplantation}\label{same-answer-content-transplantation}

A donor has the same final answer, preferentially the same proof depth,
a different orbit, and maximal proof-signature distance within the
frozen donor population:

\texttt{s\_i\ =\ L\_i(H\_donor\_same\_answer)\ -\ L\_i}.

This is a content/proof-instance specificity test. On PrOntoQA it is not
labeled a different-algorithm procedure test.

\paragraph{Different-answer
transplantation}\label{different-answer-transplantation}

This is a weaker positive control showing sensitivity to answer-bearing
task differences.

\paragraph{Semantic-twin
transplantation}\label{semantic-twin-transplantation}

The premise order is deterministically changed without changing the
protected answer, proof chain, group, or proof signature:

\texttt{t\_i\ =\ L\_i(H\_twin)\ -\ L\_i}.

A strong result requires content sensitivity together with twin
stability.

\paragraph{Truncation}\label{truncation}

The model is scored with zero through K retained latent steps to
estimate when near-clean performance appears.

\subsubsection{7. Causal status ordering}\label{causal-status-ordering}

The status is assigned lexicographically:

\begin{enumerate}
\def\labelenumi{\arabic{enumi}.}
\tightlist
\item
  replay;
\item
  necessity;
\item
  order;
\item
  structured corruption;
\item
  semantic-twin stability;
\item
  declared specificity tier.
\end{enumerate}

The first failed gate determines the terminal causal status. This
prevents a strong later statistic from compensating for an invalid
adapter or unused latent loop.

\subsubsection{8. Mode controls}\label{mode-controls}

Direct-answer mode is permitted during causal admission because it uses
no gold proof. Verbalized mode appends the protected proof steps and is
therefore inaccessible before the semantic quotient seal.

The explicit-CoT result is descriptive. It may help interpret a sealed
quotient, but it may not determine its basis, rank, thresholds, or
acceptance.

\subsubsection{9. Semantic construction
interface}\label{semantic-construction-interface}

Semantic escalation occurs only after \texttt{CONTENT\_SPECIFIC} or
\texttt{PROCEDURE\_SPECIFIC} causal admission.

The constructor sees:

\begin{itemize}
\tightlist
\item
  current latent state;
\item
  next latent state;
\item
  answer-residual future state;
\item
  anonymous query-polarity intervention response;
\item
  final answer label;
\item
  semantic-twin relation.
\end{itemize}

It does not see the protected proof-instance witness or explicit-CoT
trace.

\subsubsection{10. Candidate grammar}\label{candidate-grammar}

The finite grammar contains:

\begin{itemize}
\tightlist
\item
  answer-only rank-one direction;
\item
  PCA directions;
\item
  predictive directions;
\item
  intervention-response directions;
\item
  joint predictive-interventional directions;
\item
  equal-rank random directions;
\item
  full-state descriptive upper control.
\end{itemize}

Candidate bases are canonicalized for deterministic signs and
orthonormality.

\subsubsection{11. Candidate gates}\label{candidate-gates}

A candidate must exceed frozen thresholds for:

\begin{itemize}
\tightlist
\item
  future-state \texttt{R\^{}2};
\item
  answer-residual future-state \texttt{R\^{}2};
\item
  intervention-response \texttt{R\^{}2};
\item
  final-answer accuracy;
\item
  semantic-twin drift;
\item
  maximum rank.
\end{itemize}

Among valid candidates, the minimum rank is selected. Ties are broken by
weakest normalized margin and then a frozen method order. Full state is
not eligible.

\subsubsection{12. Semantic chronology}\label{semantic-chronology}

The order is enforced in code and tested:

\begin{enumerate}
\def\labelenumi{\arabic{enumi}.}
\tightlist
\item
  declare group-disjoint discovery, calibration, and final IDs;
\item
  capture discovery semantic orbits;
\item
  capture calibration semantic orbits;
\item
  construct candidates;
\item
  evaluate candidates only on calibration;
\item
  select quotient;
\item
  hash basis, center, split IDs, and thresholds into the quotient seal;
\item
  capture final semantic orbits only after that seal exists;
\item
  fit heads on discovery plus calibration;
\item
  evaluate final construction metrics;
\item
  reveal post-seal witness;
\item
  run causal subspace patching;
\item
  reveal explicit-CoT comparison.
\end{enumerate}

\subsubsection{13. Post-seal proof-instance
witness}\label{post-seal-proof-instance-witness}

The current PrOntoQA witness includes answer polarity, recurrent
progress, and an audit-only proof-instance coordinate derived from the
protected proof signature. It is deliberately described as
proof-instance content, not an algorithmic procedure ontology. A
stronger ProsQA witness will use explicit graph-structural features.

\subsubsection{14. Causal subspace
patching}\label{causal-subspace-patching}

For a selected basis \texttt{B}, the projector is
\texttt{P\ =\ B\ B\^{}T}. A donor semantic packet is inserted while
preserving the orthogonal complement:

\texttt{H\_patch\ =\ H\_recipient\ +\ (H\_donor\ -\ H\_recipient)\ P}.

The selected packet is compared with:

\begin{itemize}
\tightlist
\item
  full-trace transplant;
\item
  answer-only packet;
\item
  random equal-rank packet;
\item
  PCA packet as a descriptive control.
\end{itemize}

The causal gate requires a minimum recovery ratio, positive-task
fraction, and superiority over answer-only and random shortcuts.

\subsubsection{15. Statistical reporting}\label{statistical-reporting}

Effects are reported for the complete causal audit population. Bootstrap
intervals resample task rows. For final publication, proof-template or
group-level resampling should be preferred whenever public metadata
support it.

Correct and incorrect model predictions must be reported separately in
the real-result summary. A necessity-positive stratum may be reported
only as a secondary predeclared analysis; it may not replace the full
population.

\subsubsection{16. Result interpretation}\label{result-interpretation}

\begin{itemize}
\tightlist
\item
  \texttt{LATENT\_UNUSED}: no evidence that the checkpoint relies
  materially on its latent loop in this population.
\item
  \texttt{LATENT\_GENERIC}: evidence for useful ordered latent
  computation, but not enough same-answer content specificity.
\item
  \texttt{CONTENT\_SPECIFIC}: the latent trace is tied to proof-instance
  content beyond final-answer identity.
\item
  \texttt{ONTOLOGY\_UNDERIDENTIFIED}: no small quotient satisfied the
  observable causal-predictive contract.
\item
  \texttt{SEMANTIC\_CANDIDATE}: a quotient passed construction but not
  the post-seal proof-instance witness.
\item
  \texttt{SEMANTICALLY\_IDENTIFIED}: the witness passed, but subspace
  editing did not.
\item
  \texttt{CAUSALLY\_VALIDATED}: the sealed quotient passed construction,
  post-seal witness, and causal patching.
\end{itemize}

No status means that the trace is a complete or unique transcript of
hidden thought.

\subsubsection{17. Engineered positive
control}\label{engineered-positive-control}

The package includes a synthetic adapter whose trace contains a low-rank
causal content packet and high-variance nuisance dimensions. The fixture
validates the transaction and analysis mechanics. It reaches
\texttt{CAUSALLY\_VALIDATED} by design and is always marked
\texttt{engineered\_positive\_control:\ true}.

\subsubsection{18. Threats to validity}\label{threats-to-validity}

\paragraph{Adapter mismatch}\label{adapter-mismatch}

Mitigation: source-lock verification, recurrence parity test, exact
replay, fixed tokenization and suffix.

\paragraph{Final-support leakage}\label{final-support-leakage}

Mitigation: final semantic orbits are not captured until after quotient
sealing; explicit CoT is post-seal only.

\paragraph{Answer shortcut}\label{answer-shortcut}

Mitigation: same-answer donors and answer-residual future prediction.

\paragraph{Surface shortcut}\label{surface-shortcut}

Mitigation: premise-order twins and later distractor/renaming orbit
extensions.

\paragraph{Procedure overclaim}\label{procedure-overclaim}

Mitigation: PrOntoQA status is content-specific; procedure-specific
status is reserved for varied structural tasks.

\paragraph{Noise miscalibration}\label{noise-miscalibration}

Mitigation: covariance is estimated on group-disjoint calibration rows.

\paragraph{Unsupported semantic edit}\label{unsupported-semantic-edit}

Mitigation: causal patching is compared with equal-rank shortcut
controls and can fail independently of probe scores.

\paragraph{Post-result threshold
changes}\label{post-result-threshold-changes}

Mitigation: config, model, source, dataset, donors, splits, and output
contract are sealed before execution.

\subsubsection{19. Hardware plan}\label{hardware-plan}

The primary target is intentionally small. The recommended configuration
for an 8 GB RTX 3070 Ti is batch size one, deterministic inference,
short prompts, capture caching, and moving saved traces to CPU. The
package does not retrain the external model.

\subsubsection{20. Escalation path}\label{escalation-path}

\begin{enumerate}
\def\labelenumi{\arabic{enumi}.}
\tightlist
\item
  Run 12-task smoke solely to validate source parity and replay.
\item
  Run 64-task primary without changing gates.
\item
  Inspect the canonical status and failure coordinate.
\item
  Run the 256-task profile only if the adapter remains valid.
\item
  Add ProsQA only after PrOntoQA produces a stable causal result.
\item
  Add CODI or another latent method for cross-model transport only after
  one checkpoint has a valid semantic object.
\item
  Add SLPO for adaptive-depth analysis.
\item
  Add Hyper as a scale stress.
\end{enumerate}

\subsubsection{21. Stopping rules}\label{stopping-rules}

\begin{itemize}
\tightlist
\item
  Replay failure: stop and repair the adapter; do not interpret latent
  interventions.
\item
  Out-of-memory during smoke: reduce runtime memory, not scientific
  thresholds.
\item
  \texttt{LATENT\_UNUSED}: preserve as the primary result; do not filter
  to successful tasks as the headline.
\item
  \texttt{LATENT\_GENERIC}: do not run semantic escalation.
\item
  No valid quotient: report \texttt{ONTOLOGY\_UNDERIDENTIFIED}.
\item
  Causal patch failure: do not call the quotient causally validated.
\end{itemize}

\subsubsection{22. Current execution
status}\label{current-execution-status}

The source-locked external checkpoint was not executed in the
environment used to build this package. That environment had CPU-only
PyTorch, no Transformers installation, no external checkout, no
checkpoint, no public dataset, and no CUDA device. The package's
validated result is the engineered fixture only.

\subsubsection{23. Decision after the first real
receipt}\label{decision-after-the-first-real-receipt}

The first real receipt should determine the next move:

\begin{itemize}
\tightlist
\item
  adapter invalid -\textgreater{} repair only adapter fidelity;
\item
  latent unused -\textgreater{} test a second published checkpoint or
  task, without rewriting the result;
\item
  latent generic -\textgreater{} compare latent and direct modes and
  stop semantic claims;
\item
  content specific -\textgreater{} execute sealed quotient construction;
\item
  semantic candidate -\textgreater{} improve independent witness design
  in a new protocol version;
\item
  causally validated -\textgreater{} proceed to ProsQA structure
  sensitivity and cross-model transport.
\end{itemize}

\newpage

\section{Engineered positive-control
result}\label{engineered-positive-control-result}

The included fixture is deterministic and deliberately contains a
low-rank causal semantic packet. It validates the mechanics only; it is
not a trained language model.

\begin{longtable}[]{@{}lr@{}}
\toprule\noalign{}
Quantity & Frozen result \\
\midrule\noalign{}
\endhead
\bottomrule\noalign{}
\endlastfoot
Fixture tasks & 48 \\
Covariance calibration tasks & 8 \\
Causal audit tasks & 40 \\
Causal status & \texttt{CONTENT\_SPECIFIC} \\
Semantic status & \texttt{CAUSALLY\_VALIDATED} \\
Selected quotient & \texttt{joint}, rank 3 \\
Replay maximum & 0 \\
Necessity & 5.814125 {[}5.655045, 5.964764{]} \\
Order & 0.829070 {[}0.784764, 0.874355{]} \\
Norm-matched structure & 7.342361 {[}7.027469, 7.659845{]} \\
Covariance-matched structure & 12.649418 {[}12.020959, 13.362811{]} \\
Same-answer content transplant & 4.408597 {[}3.409079, 5.523186{]} \\
Different-answer transplant & 14.829357 {[}13.318282, 16.259007{]} \\
Semantic-twin transplant & 0.000000 {[}0.000000, 0.000000{]} \\
Final future R2 & 0.751820 \\
Final answer-residual future R2 & 0.636851 \\
Final intervention R2 & 0.984125 \\
Post-seal proof-witness R2 & 0.954829 \\
Causal patch passed & \texttt{true} \\
\end{longtable}

\begin{Shaded}
\begin{Highlighting}[]
\NormalTok{fixture pre{-}audit seal: 68b04ea973f5030a7085d2b1891e87bedbf058db007d004fefee047beaee70b9}
\NormalTok{fixture result receipt: 924157159d14b8d3715f2f886c2ad294b55f5cac15a3b63f1552df1b3731e620}
\NormalTok{independent verification: e4e0c5806022faa79e41a363a2100695a8b9d0f33a85c931014f8c732b84c9ca}
\end{Highlighting}
\end{Shaded}

This result validates the decision logic, evidence chronology, and
intervention plumbing. It establishes nothing about the public
checkpoint.

\section{Current runtime boundary}\label{current-runtime-boundary}

\begin{Shaded}
\begin{Highlighting}[]
\FunctionTok{\{}
  \DataTypeTok{"checkpoint\_exists"}\FunctionTok{:} \KeywordTok{false}\FunctionTok{,}
  \DataTypeTok{"cuda\_available"}\FunctionTok{:} \KeywordTok{false}\FunctionTok{,}
  \DataTypeTok{"data\_exists"}\FunctionTok{:} \KeywordTok{false}\FunctionTok{,}
  \DataTypeTok{"external\_repo\_exists"}\FunctionTok{:} \KeywordTok{false}\FunctionTok{,}
  \DataTypeTok{"ready"}\FunctionTok{:} \KeywordTok{false}\FunctionTok{,}
  \DataTypeTok{"torch"}\FunctionTok{:} \StringTok{"2.10.0+cpu"}\FunctionTok{,}
  \DataTypeTok{"transformers\_error"}\FunctionTok{:} \StringTok{"No module named \textquotesingle{}transformers\textquotesingle{}"}
\FunctionTok{\}}
\end{Highlighting}
\end{Shaded}

The third-party checkpoint result remains unknown.

\section{Part III - Source interface and
chronology}\label{part-iii---source-interface-and-chronology}

The adapter is valid only when it satisfies all of the following.

\subsection{Trace identity}\label{trace-identity}

\texttt{capture(task)} returns the vectors actually reinjected at the
model's latent-token positions, in recurrent order, with shape
\texttt{{[}latent\_steps,\ hidden\_size{]}}.

\subsection{Supplied-trace replay}\label{supplied-trace-replay}

\texttt{score\_with\_trace(task,\ trace)} runs the same teacher-forced
prompt and answer suffix while replacing only the reinjected latent
vectors. Exact own-trace replay must reproduce answer logits within the
frozen tolerance.

\subsection{Source parity}\label{source-parity}

The package includes a separate reference recurrence derived from the
pinned upstream source. The adapter engine must match it exactly on a
controlled PyTorch fixture.

\subsection{Fixed surfaces}\label{fixed-surfaces}

Across clean and intervention arms, the following remain fixed:

\begin{itemize}
\tightlist
\item
  checkpoint and base model;
\item
  tokenizer and special tokens;
\item
  prompt text for the recipient task;
\item
  token positions;
\item
  attention mask;
\item
  answer suffix;
\item
  teacher-forced label mask;
\item
  decoding settings.
\end{itemize}

\subsection{Modes}\label{modes}

\begin{itemize}
\tightlist
\item
  \texttt{latent}: admissible for causal and semantic construction.
\item
  \texttt{direct}: admissible as a pre-seal mode control.
\item
  \texttt{verbalized}: contains gold proof steps and is post-seal only.
\end{itemize}

\subsection{Fingerprint}\label{fingerprint}

The adapter fingerprint commits:

\begin{itemize}
\tightlist
\item
  upstream repository and commit;
\item
  source-verification receipt;
\item
  checkpoint repository, revision, filename, and full SHA-256;
\item
  local base-model path, revision, and tree SHA-256;
\item
  latent-step count;
\item
  device type.
\end{itemize}

\subsection{Failure}\label{failure}

Any replay, source-lock, cache representation, latent-position, or
trace-shape mismatch returns an adapter failure. It must not be treated
as evidence about the model's semantics.

\begin{Shaded}
\begin{Highlighting}[]
\NormalTok{T0  package code, thresholds, claim tier, source lock frozen}
\NormalTok{T1  external repository checked out at exact commit}
\NormalTok{T2  required Git blobs verified and dirty tree rejected}
\NormalTok{T3  pinned checkpoint and GPT{-}2 revision downloaded}
\NormalTok{T4  full checkpoint and local GPT{-}2 payload hashed}
\NormalTok{T5  deterministic task population and group IDs frozen}
\NormalTok{T6  covariance{-}calibration and causal{-}audit groups frozen}
\NormalTok{T7  donor plans and semantic split IDs frozen}
\NormalTok{T8  pre{-}audit seal written}
\NormalTok{T9  calibration traces captured and covariance committed}
\NormalTok{T10 causal audit executed on disjoint audit rows}
\NormalTok{T11 causal status assigned}
\NormalTok{T12 if admitted, semantic discovery traces captured}
\NormalTok{T13 semantic calibration traces captured}
\NormalTok{T14 candidate quotient chosen}
\NormalTok{T15 quotient basis, center, thresholds, and splits sealed}
\NormalTok{T16 semantic final traces captured for the first time}
\NormalTok{T17 post{-}seal proof{-}instance witness revealed}
\NormalTok{T18 causal subspace patching executed}
\NormalTok{T19 explicit{-}CoT comparison revealed}
\NormalTok{T20 canonical result receipt written}
\end{Highlighting}
\end{Shaded}

No T16-T19 object may influence T14. The implementation records and
tests the semantic chronology.

The external assets are intentionally not redistributed.

\subsection{Source repository}\label{source-repository}

\begin{itemize}
\tightlist
\item
  \texttt{connordilgren/are-lrms-easily-interpretable}
\item
  pinned commit: \texttt{32f413d8d55239d9bc54bb6b6ec37b0630891ed4}
\end{itemize}

The package verifies the wrapper, latent runtime, requirements, and
selected dataset Git blobs listed in
\texttt{external\_source\_lock.json}.

\subsection{Primary checkpoint}\label{primary-checkpoint}

\begin{itemize}
\tightlist
\item
  repository: \texttt{connordilgren/gpt2-prontoqa-multimode-coconut}
\item
  revision: \texttt{4ec759a9f9b8a8e7e1a8d1e3e2d43eef969ac83f}
\item
  file: \texttt{checkpoint\_37}
\end{itemize}

\subsection{Base model}\label{base-model}

\begin{itemize}
\tightlist
\item
  repository: \texttt{openai-community/gpt2}
\item
  revision: \texttt{607a30d783dfa663caf39e06633721c8d4cfcd7e}
\end{itemize}

\subsection{Runtime}\label{runtime}

The isolated real environment pins PyTorch 2.5.1 and Transformers
4.46.2, matching the reviewed external requirements where exact pins
were available.

\subsection{Secondary planned
checkpoint}\label{secondary-planned-checkpoint}

ProsQA source and checkpoint revisions are recorded in
\texttt{external\_source\_lock.json}. They are not executed by the
primary v0.3.0 run path.

\newpage

\section{Part IV - Real execution status and
runbook}\label{part-iv---real-execution-status-and-runbook}

\subsection{Canonical status}\label{canonical-status}

\texttt{EXTERNAL\_EXECUTION\_NOT\_RUN}

The build environment used for v0.3.0 reported:

\begin{itemize}
\tightlist
\item
  CPU-only PyTorch available;
\item
  CUDA unavailable;
\item
  Transformers unavailable;
\item
  pinned upstream checkout absent;
\item
  checkpoint absent;
\item
  public dataset absent.
\end{itemize}

No attempt was made to substitute a different model or fabricate
external metrics. The external checkpoint remains an open empirical
test.

\subsection{What was executed}\label{what-was-executed}

\begin{itemize}
\tightlist
\item
  engineered fixture generation;
\item
  engineered pre-audit sealing;
\item
  causal and semantic pipeline on the fixture;
\item
  automated test suite;
\item
  source-lock unit tests;
\item
  recurrence parity fixture;
\item
  chronology and split-isolation tests;
\item
  independent recomputation of fixture means and semantic-basis
  properties;
\item
  package closure and manifest verification.
\end{itemize}

\subsection{What must happen on the user's
machine}\label{what-must-happen-on-the-users-machine}

\begin{enumerate}
\def\labelenumi{\arabic{enumi}.}
\tightlist
\item
  Run \texttt{Setup\_External\_PrOntoQA\_Coconut.ps1}.
\item
  Run the smoke profile.
\item
  Inspect replay and source verification.
\item
  Run the primary profile without modifying the frozen config.
\item
  Preserve the complete \texttt{frozen}, \texttt{results},
  \texttt{capture\_cache}, and environment records.
\end{enumerate}

\subsection{Failure policy}\label{failure-policy}

The first real status is the result. An adapter failure permits an
adapter-only revision because no semantic claim is yet admissible. Any
scientific threshold revision requires a new protocol and new final
support.

\subsection{Before setup}\label{before-setup}

\begin{itemize}
\tightlist
\item[$\square$]
  At least 25 GB free disk space.
\item[$\square$]
  NVIDIA driver operational.
\item[$\square$]
  Python 3.10, 3.11, or 3.12 available through \texttt{py}.
\item[$\square$]
  Git available in PowerShell.
\item[$\square$]
  Network access to GitHub, PyPI, PyTorch wheels, and Hugging Face.
\end{itemize}

\subsection{After setup}\label{after-setup}

\begin{itemize}
\tightlist
\item[$\square$]
  \texttt{runtime\_paths.json} exists.
\item[$\square$]
  upstream commit equals the source lock.
\item[$\square$]
  source blob verification passes.
\item[$\square$]
  checkpoint revision and full SHA-256 are recorded.
\item[$\square$]
  GPT-2 revision and local tree SHA-256 are recorded.
\item[$\square$]
  \texttt{pip\ check} passes.
\item[$\square$]
  CUDA is visible to the isolated environment.
\end{itemize}

\subsection{Smoke run}\label{smoke-run}

\begin{itemize}
\tightlist
\item[$\square$]
  Seal is written before model audit.
\item[$\square$]
  Exact replay passes all 12 rows.
\item[$\square$]
  Prompt, suffix, latent positions, and answer mask are stable.
\item[$\square$]
  No out-of-memory or cache-format error occurs.
\end{itemize}

\subsection{Primary run}\label{primary-run}

\begin{itemize}
\tightlist
\item[$\square$]
  Configuration file unchanged from the sealed copy.
\item[$\square$]
  Covariance calibration rows are group-disjoint from audit rows.
\item[$\square$]
  Donor plans cover every causal audit row once.
\item[$\square$]
  Explicit CoT remains unavailable before semantic seal.
\item[$\square$]
  Final semantic orbits are captured only after quotient seal.
\item[$\square$]
  Complete result and environment records are retained even on failure.
\end{itemize}

\subsection{After result}\label{after-result}

\begin{itemize}
\tightlist
\item[$\square$]
  Run \texttt{verify-result}.
\item[$\square$]
  Generate summary.
\item[$\square$]
  Copy the complete result directory before any exploratory analysis.
\item[$\square$]
  Do not relabel content specificity as algorithmic procedure
  specificity.
\end{itemize}

\subsection{Objective}\label{objective}

Execute a source-locked causal-semantic audit of the public multimode
GPT-2 Coconut PrOntoQA checkpoint without changing the v0.3.0 gates
after result access.

\subsection{Recommended first command}\label{recommended-first-command}

\begin{Shaded}
\begin{Highlighting}[]
\OperatorTok{.}\NormalTok{\textbackslash{}Setup\_External\_PrOntoQA\_Coconut}\OperatorTok{.}\FunctionTok{ps1}
\end{Highlighting}
\end{Shaded}

Then:

\begin{Shaded}
\begin{Highlighting}[]
\OperatorTok{.}\NormalTok{\textbackslash{}Run\_Real\_LRM\_Pilot}\OperatorTok{.}\FunctionTok{ps1} \OperatorTok{{-}}\NormalTok{Profile smoke }\OperatorTok{{-}}\NormalTok{Device cuda}
\end{Highlighting}
\end{Shaded}

The smoke run is an adapter and replay qualification. It has semantic
escalation disabled.

\subsection{Advancement rule}\label{advancement-rule}

Proceed to \texttt{primary} only when:

\begin{itemize}
\tightlist
\item
  source verification passes;
\item
  the checkpoint and base-model hashes are recorded;
\item
  exact replay is below tolerance;
\item
  no cache or token-position discrepancy is observed;
\item
  GPU memory is stable.
\end{itemize}

\subsection{Scientific rule}\label{scientific-rule}

Do not alter \texttt{configs/prontoqa\_multimode\_coconut\_gpt2.json}
after staging. If the primary result is negative, preserve it. A new
hypothesis requires a new version and seal.

\subsection{Returned files to
preserve}\label{returned-files-to-preserve}

\begin{itemize}
\tightlist
\item
  \texttt{external\_runtime/environment.json}
\item
  \texttt{external\_runtime/runtime\_paths.json}
\item
  \texttt{frozen/real\_primary\_preaudit\_seal.json}
\item
  donor plans and experiment plan beside the seal
\item
  \texttt{results/real\_primary\_block24.json}
\item
  \texttt{results/real\_primary\_summary.md}
\item
  \texttt{results/capture\_cache/} if present
\end{itemize}

\subsection{Current package evidence}\label{current-package-evidence}

The engineered positive control passed. The external model has not been
executed. No real-model conclusion should be inferred from the fixture.

\newpage

\section{Part V - Structure-sensitive
follow-up}\label{part-v---structure-sensitive-follow-up}

\subsection{Purpose}\label{purpose-1}

PrOntoQA is the source-faithful causal admission tier. ProsQA is the
intended structure-sensitive tier. It should be activated only after the
PrOntoQA adapter passes replay.

\subsection{Stronger claim}\label{stronger-claim}

A ProsQA result may use \texttt{PROCEDURE\_SPECIFIC} only when all of
the following are implemented before final audit:

\begin{enumerate}
\def\labelenumi{\arabic{enumi}.}
\item
  proof graphs or derivation structures are extracted independently of
  text surface;
\item
  donors are matched on final answer and broad difficulty;
\item
  a structure-near donor and structure-far donor are assigned for every
  recipient;
\item
  the primary statistic is the contrast

  \texttt{loss(far\ donor)\ -\ loss(near\ donor)};
\item
  the contrast has a positive frozen interval and positive-task
  fraction;
\item
  semantic twins preserve proof topology while changing wording, entity
  names, and premise order;
\item
  final proof structures are held out by template family.
\end{enumerate}

\subsection{Why the contrast matters}\label{why-the-contrast-matters}

A harmful arbitrary donor can reflect task identity, vocabulary, or
context formatting. The far-minus-near contrast asks whether disruption
grows with formal solution-structure distance after answer and
difficulty are controlled.

\subsection{Additional adapter work}\label{additional-adapter-work}

ProsQA answers are not limited to Boolean labels. A separate answer
parser and protected loss contract are required. The current PrOntoQA
adapter must not be silently reused as if these formats were identical.

\subsection{Frozen source references}\label{frozen-source-references}

The external source lock records the ProsQA test-data Git blob,
checkpoint repository, checkpoint revision, and checkpoint filename.
Those entries are planning locks, not evidence that the ProsQA
checkpoint has been executed.

\section{Part VI - Claim boundary and decision
policy}\label{part-vi---claim-boundary-and-decision-policy}

\subsection{Current release status}\label{current-release-status}

\texttt{INFRASTRUCTURE\_VALIDATED\_EXTERNAL\_EXECUTION\_PENDING}

The package validates a source-faithful intervention architecture, a
staged chronology, and an engineered causal-semantic positive control.
The real public checkpoint was not run in the build environment.

\subsection{Supported claims}\label{supported-claims}

The release supports the following statements:

\begin{enumerate}
\def\labelenumi{\arabic{enumi}.}
\tightlist
\item
  The adapter contains a separately tested implementation of the
  reviewed Coconut recurrence and exposes supplied-trace replay.
\item
  Exact source commit and required Git blobs are pinned and verifiable.
\item
  Checkpoint and base-model revisions are pinned by setup; local payload
  hashes are sealed at staging.
\item
  Covariance-matched noise uses a group-disjoint calibration support in
  the Block 24 pipeline.
\item
  Semantic discovery, calibration, quotient sealing, final capture,
  proof witness, causal patching, and explicit-CoT comparison occur in
  the declared order.
\item
  The engineered fixture reaches \texttt{CAUSALLY\_VALIDATED} and all
  included automated tests pass.
\item
  PrOntoQA's primary maximum specificity claim is content/proof-instance
  specificity, not proof of multiple internal reasoning algorithms.
\end{enumerate}

\subsection{Unsupported claims}\label{unsupported-claims}

The release does not establish that the real Coconut checkpoint:

\begin{itemize}
\tightlist
\item
  relies on latent recurrence;
\item
  carries ordered latent computation;
\item
  represents proof-instance content beyond answer identity;
\item
  is stable under semantic twins;
\item
  exposes a low-rank semantic quotient;
\item
  supports causal semantic editing;
\item
  matches its explicit natural-language proof mode;
\item
  generalizes to ProsQA or any other task;
\item
  is safe, conscious, sentient, or recursively self-improving.
\end{itemize}

The engineered fixture cannot be used as evidence for any external
model.

\subsection{Interpretation boundary}\label{interpretation-boundary}

Even a real \texttt{CAUSALLY\_VALIDATED} result would mean only that,
within the pinned model, task population, intervention language, and
frozen thresholds, a compact latent subspace predicts and causally
mediates declared proof-instance consequences. It would not reveal a
verbatim chain of thought or prove that the discovered coordinates are
the unique correct semantics.

\subsection{Revision boundary}\label{revision-boundary}

Any threshold, donor rule, counterfactual generator, source revision,
checkpoint, model mode, or claim tier changed after a real result must
produce a new protocol version and a new pre-audit seal. A failed run
remains part of the evidence lineage.

The next action is the 12-task source-parity and replay smoke run on the
user's RTX 3070 Ti.

After the primary PrOntoQA result:

\begin{itemize}
\tightlist
\item
  \texttt{ADAPTER\_INVALID}: repair adapter fidelity only.
\item
  \texttt{LATENT\_UNUSED}: preserve the result and consider a different
  public latent checkpoint.
\item
  \texttt{LATENT\_GENERIC}: stop semantic escalation; compare execution
  modes descriptively.
\item
  \texttt{CONTENT\_SPECIFIC}: complete the sealed semantic escalation.
\item
  \texttt{CAUSALLY\_VALIDATED}: proceed to the ProsQA
  structure-sensitive protocol before any broader claim.
\end{itemize}

Default recommendation: do not proceed directly to Hyper or training.
Resolve the smallest external checkpoint first.

\end{document}
